\documentclass[a4paper]{article}

\usepackage{anyfontsize}
\usepackage{amsmath}

\usepackage[no-math]{fontspec}
\setmainfont{TeX Gyre Termes}
\usepackage{unicode-math}
\unimathsetup{bold-style=ISO,sans-style=italic}
\setmathfont{XITS Math}
\AtBeginDocument{
  \let\uglymathbf\mathbf
  \renewcommand\mathbf\symbf
  \let\uglymathsf\mathsf
  \renewcommand\mathsf\symsf
}

\AtBeginDocument{
  \let\uglyepsilon\epsilon
  \renewcommand\epsilon\varepsilon
}

\usepackage{indentfirst}
\setlength{\parindent}{0em}
\usepackage{autobreak}
\allowdisplaybreaks
\usepackage{parskip}
\usepackage{geometry}
\geometry{top=3cm,bottom=3cm,left=2cm,right=2cm}

\usepackage[fontset=none]{ctex}
\setCJKmainfont{Adobe Song Std}[BoldFont=Noto Sans CJK SC Medium]

\begin{document}

\title{\textbf{实测天体物理作业}}
\author{1900011604\ \ 张植竣\ \ 北京大学}
\date{2022年2月21日}
\maketitle

\begin{enumerate}

    \item[1.]
    使用Rayleigh-Jeans近似：
    \[
        B_\nu = \frac{2\nu^2kT}{c^2} = 4.0094\times10^{-18}\,\mathrm{W/Hz}
    \]
    太阳对观测者所张的立体角为：
    \[
        \Omega = \frac{\pi R^2}{r^2} = 6.7919\times10^{-5}\,\mathrm{sr}
    \]
    则接收到的流量密度为：
    \[
        S_\nu = B_\nu\Omega = 2.723\times10^{-22}\,\mathrm{W/(m^2\cdot Hz)} = 2.723\times10^4\,\mathrm{Jy}
    \]

    \item[2.]
    使用Rayleigh-Jeans近似：
    \[
        B_\nu = \frac{2\nu^2kT}{c^2} = 4.0094\times10^{-18}\,\mathrm{W/Hz}
    \]
    太阳对观测者所张的立体角为：
    \[
        \Omega = \frac{\pi R^2}{r^2} = 1.5965\times10^{-15}\,\mathrm{sr}
    \]
    则接收到的流量密度为：
    \[
        S_\nu = B_\nu\Omega = 6.4\times10^{-33}\,\mathrm{W/(m^2\cdot Hz)} = 0.64\,\mu\mathrm{Jy}
    \]
    因为$0.64\,\mu\mathrm{Jy}<10\,\mu\mathrm{Jy}$，无法探测。

\end{enumerate}

\end{document}